\smallsec{Vision-based Object Manipulation} The learning of robotic control policy has been traditionally studied with low-dimensional state observations \cite{andrychowicz2020learning, hwangbo2019learning, kumar2021rma, schulman2017proximal, xu2019learning}. Recently, vision-based policies~\cite{anderson2018evaluation,goyal2020packit,shah2021ving,yang2021learning,huang2021generalization,chen2021system,yu2021visuallocomotion,shah2022offline} have gained increasing attention since the high-dimensional visual sensory input provides more generalizable observation representation across tasks and is more accessible in real-world perception systems. Various forms of visual input have been explored. Prior work has directly encoded the RGB images into a low-dimensional latent space and relied on model-based \cite{hafner2019dream, watter2015embed} or model-free \cite{yarats2021image, yarats2022mastering} reinforcement learning (RL) to train policies to operate in this space. More recently, RT-1 \cite{rt12022arxiv} infers the robot's actions from a history of images by leveraging transformer architectures ~\cite{vaswani2017attention}. Our proposed \method also uses a transformer to predict actions, however, unlike RT-1, we additionally leverage depth to construct a multi-view scene representation. The use of depth input has also been extensively studied. Methods like CLIPort \cite{shridhar2021cliport} and IFOR \cite{goyal2022ifor} directly process the RGB-D images for object manipulation, and hence are limited to simple pick-and-place tasks in 2D top-down settings. To overcome this issue, explicit 3D representations such as point clouds have been utilized. C2F-ARM \cite{c2farm} and PerAct \cite{peract2022arxiv} voxelize the point clouds and use a 3D convolutional network as the backbone for control inference. However, high-precision tasks typically require high resolution of voxelization, resulting in high memory consumption and slow training. Our approach falls into this category but addresses the scalability issue by transforming the point cloud into a set of RGB-D images from multiple views. We show that this significantly improves memory footprint and training efficiency, and leads to higher performance when compared to directly working with RGB(-D) or point cloud input (see Table.~\ref{table:rlbench}). Another relevant work is MIRA \cite{lin2022mira}, which also uses novel view images to represent the 3D scene for action inference. MIRA achieves this by implicitly constructing a neural radiance field representation (NeRF) of the scene from a set of RGB images and then generating novel view images from the optimized NeRF model. However, the requirement of optimizing a scene NeRF model slows down the inference speed at test time and relies on RGB images from a dense set of views as input. In contrast, our approach can achieve significantly faster inference speed and can work with even a single-view RGB image.

\smallsec{Multi-Task Learning in Robotics} Learning a single model for many different tasks has been of particular interest to the robotics community recently. A large volume of work achieves the multi-task generalization by using a generalizable task or action representation such as object point cloud \cite{huang2021generalization, chen2021system}, semantic segmentation and optical flow \cite{goyal2022ifor}, and object-centric representation \cite{lenz2015deep, saxena2006robotic}. However, the limited expressiveness of such representations constrains them to only generalize within a task category. Task parameterization \cite{deisenroth2014multi, kalashnikov2021mt} and discrete task-dependent output layer \cite{devin2017learning, teh2017distral} approaches are investigated with reinforcement learning to learn policies for tasks in different categories. With the recent breakthrough in large language models, multi-task robot learning has been approached by using natural language to specify a broad range of tasks and learning the policy from large pre-collected datasets \cite{saycan2022arxiv, rt12022arxiv, goyal2020rel3d, jang2021bcz, lynch2020language, mees2022matters, nair2022learning, reed2022generalist, singh2022progprompt}. We are inspired by this success but propose to learn language-conditioned multi-task policies with a small demonstration dataset.

\smallsec{Transformers for Object Manipulation} The success of transformers in vision and NLP has led its way into robot learning \citep{chaplot2021differentiable, clever2021assistive,johnson2021motion,yang2021learning}. Especially in object manipulation, transformer-based models with an attention mechanism can be utilized to extract features from sensory inputs to improve policy learning~\cite{cachet2020transformer,kim2021transformer,dasari2021transformers,jangir2022look,liu2022structformer}. Unlike most prior work, we do not use large datasets for training. \method efficiently learns from a small set of demonstrations, handle multiple views as visual inputs, and fuses information from language goals to tackle multiple manipulation tasks.

\smallsec{Multi-View Networks in Computer Vision} Multi-view representations have been explored in various vision problems. For point cloud recognition, SimpleView~\cite{goyal2021revisiting} showed how a simple view-based method outperforms sophisticated point-based methods. Follow-up works like MVTN~\cite{hamdi2021mvtn} and Voint cloud~\cite{hamdi2021voint} have further improved upon SimpleView's architecture. Multi-view representations have also been used for other problems like 3D visual grounding~\cite{huang2022multi} and view synthesis~\cite{mildenhall2021nerf}. Unlike them, we focus on the problem of predicting robot actions for object manipulation.